\pg{426} \textsc{Deus} demonstra-se pela existencia das creaturas; porque estas, por serem \emph{effeitos}, suppoóem uma causa proporcionada, que é \textsc{Deus}. Logo a existencia de \textsc{Deus} demonstra-se \emph{a-posteriori} \footnote{Cf. S. Thom., --- \emph{Sum}. \emph{Th}., p. I, q. 2, a. 2; --- \emph{c}. \emph{Gentes}, l. I, c. 12.\par A demonstração da existencia de Deus basea-se no \emph{principio de causalidade}. E' por isso que os \emph{materialistas}, os \emph{positivistas} e todos os \emph{incredulos}, no intuito de negar esta verdade fundamental, ou, ao menos, de prescindir d'ella, atacam a noção de causa e esse principio de \emph{causalidade}. --- Não repetimos tudo o que na \emph{Ontologia} escrevemos em defeza do dito principio. Limitamo-nos a observar que o principio de \emph{causalidade} é uma verdade, que a intelligencia descobre e formula á vista das mudanças, dos phenomenos, dos factos, que se succedem, uns aos outros, no mundo visivel. Ora quando um phenomeno começa, percebe-se immediatamente que está relacionado com um outro, de que deriva. Esta relação não é apenas de \emph{successdo}, mas de \emph{dependencia}. Tal dependencia é uma coiea real, objectiva, existente nos proprios seres creados; a intelligencia \emph{descobre-a}, mas \emph{não a inventa}. Por isso o principio de \emph{causalidade} tem un valor real e objectivo; porque, embora este juizo --- \emph{todo o effeito tem a sua causa} --- seja um producto da razão, comtudo a rasdo \emph{ndo o produz senão fundada nos dados da experiencia sensivel}. A experiencia apresenta os factos; a razão descobre o nexo, que os liga. ---\pgn{427}Advirta-se que o principio de \emph{causalidade} não se applica a todo o ser que \emph{existe}, mas a todo o ser que começa a existir, e por isso a Causa \emph{primeira} não tem causa.\par Os \emph{Kantistas} não negam o principio de \emph{causalidade}, mas dizem que elle só tem valor, quando se trata de dois phenomenos, ambos sujeitos á observação sensivel, porque, em tal caso, podemos averiguar a dependencia de um phenomeno do outro, --- mas não tem valor, quando se trata de duas realidades, uma das quaes não está sujeita á experiencia, porque, em tal caso, não podemos conhecer a \emph{dependencia} ou a \emph{origem} de uma realidade da outra, e por isso nem a razão de \emph{causa} e de \emph{ejfeíto}; o que se verifica no nosso caso, porque, se o mundo está sujeito á experiencia, Deus escapa a toda a experiencia. Por isso não podemos basear-nos no principio de \emph{causalidade}, para deduzir da existencia do mundo a existencia de Deus.\par Respondemos que, para o principio de \emph{causalidade} ter todo o seu valor, não é necessario que ambas as realidades, que se dizem proceder uma da outra, estejam sujeitas á experiencia sensivel. Basta que um phenomeno, um facto, presente aos nossos sentidos, manifeste um caracter de mutabilidade e de contingencia, para que possamos e devamos concluir que esse phenomeno \emph{depende} necessariamente de um antecedente. Pouco importa, para O caso, que esse antecedente seja visivel ou invisivel, natural ou sobrenatural; elle existe, deve existir, porque existe o phenomeno, que não tem em si a razão sufficiente da sua existencia. Ninguem viu a força da \emph{gravitação}; e, comtudo, todos a admittem, porque se revela na queda dos graves. Ninguem viuo talento militar de Napoleão; mas todos o admittem, porque se manifestou nos seus grandes planos de guerra. Deus é invisivel na sua essencia, mas revela-se nas coisas creadas, que suppõem e exigem uma força infinita. Portanto, ou devemos negar o principio de \emph{causalidade}, e temos a completa derrocada de todas as sciencias, --- ou devemos admittir que, em força d'esse principio, podemos e devemos elevar-nos ao conhecimento da existencia de Deus. --- Nem se diga que uma percepção mais profunda e completa dos phenomenos naturaes pode fazer com que se prescinda de uma causa primeira. Os phenomenos naturaes, ainda que podessem ser percebidos em toda a sua comprehensão, nunca podem perder os \emph{essenciaes} caracteres de \emph{contingencia} e de \emph{mutabilidade}, e por isso da sua \emph{dependencia} de uma causa primeira. Pelo contrario: quanto mais se conhece o \emph{movimento}, a \emph{larmonia}, a \emph{finalidade} dos seres visiveis, tanto mais necessaria apparece a «xistencia de um supremo e omnipotente Ordenador.\par Objectam: Todos os nossos conhecimentos são \emph{relativos} ao mundo \emph{visive}? o á intelligencia \emph{finita} do homem. Logo não podemos conhecer a \textsc{Deus}, que se suppõe \emph{invisivel} e \emph{infinito}.\par Resposta. E' certo que todos os nossos conhecimentos, \emph{que adquirimos}\pgn{428}\emph{de um modo directo e immediato}, se referem ao mundo visivel. Mas d'ahi não se segue que a nossa razão não possa, nem de um modo \emph{indirecto} e \emph{mediato}, adquirir outros conhecimentos, relativos ao mundo invisivel. Os caracteres de contingencia e de mutabilidade, que observamos no mundo visivel, levam-nos necessariamente a admittir a existencia de um Ente invisivel, necessario e immutavel. --- Como tambem é certo que todos os conhecimentos são proporcionados á capacidade da intelligencia cognoscente e ao modo por que os objectos se tornam presentes á mesma intelligencia, de modo que só apprehendemos as notas do objecto, que actuaram na faculdade, ficando incognosciveis as que escaparam ao alcance d'ella. Mas esta \emph{relatividade}, se prova que não podemos \emph{comprehender} \& Deus, porque, a intelligibilidade de Deus excede infinitamente a capacidade da intelligencia creada, não prova que não possamos, \emph{de algum modo}, conhecer o omnipotente Creador de todas as coisas.\par .Devemos, porém, advertir que, para os sequazes do \emph{agnosticismo}, \emph{relatividade} significa \emph{subjectivismo}; e por isso quando dizem que os nossos conhecimentos são \emph{relativos}, entendem que são \emph{subjeclivos}, em quanto nos fazem conhecer as coisas não como são em si mesmas, mas como são no nosso espirito, ás quaes nada corresponde na realidade. Não nos detemos na refutação d'este erro, a qual já está feita. Só acerescentaremos que, se todos os nossos conhecimentos sao \emph{swbjectivos}, são \emph{subjectivos} não só os que se referem ao mundo invisivel, mas tambem 03s que se referem ao visivel, e por isso, se não existe Deus, nem existe coisa alguma no mundo. --- Nem se diga que os conhecimentos relativos ao mundo \emph{visivel} são \emph{objectivos}, porque podem ser verificados pela experiencia, pelos factos: o que se não pode fazer com os conhecimentos relativos ao mundo \emph{invisivel}, a Deus. Negamos que os nossos conhecimentos relativos a Deus não se possam verificar. pelos factos. Na verdade, como provamos nós que Deus existe? Provamol-o pela origem e natureza do mundo, do movimento, da vida, da ordem, da finalidade; porque a origem e a natureza d'estas coisas suppõem a existencia de uma causa extramundana. Ora o mundo, o movimento, a vida, a ordem, a finalidade, são factos, e factos grandiosos, admiraveis, que actuam nos nossos sentidos e despertam a admiração na nossa intelligencia. Portanto o conhecimento . de Deus é tão \emph{objectivo}, como \emph{objectivos} são os factos, que nos servem de degrau para nos elevarmos a Elle.\par Continuam: A existencia de Deus não pode ser demonstrada pelos effeitos. Na verdade, entre a causa e o effeito deve existir uma proporção; ora nenhuma proporção pode existir entre as creaturas, que são seres finitos, e Deus, que é o Ser infinito.\par Resposta. Demonstra-se á causa pelos effeitos por dois modos. Se os effeitos possuem a mesma natureza da causa (0 que se verifica nas causas \emph{univocas}) então não só levam ao conhecimento da \emph{existencia}, mas tambem da \emph{natureza} da cansa, de modo que adquirimos um conhecimento \emph{perfeito} da causa. Se os effeitos não têm a mesma natureza da causa (\emph{o} que se verifica nas causas \emph{equivocas}), então só leyam ao conhecimento da \emph{ewístencia} da causa, e não da \emph{natureza}. e por isso o conhecimento da causa é \emph{imperfeito}\pgn{429}, e o grau d'esta imperfeição é proporcionado ao grau de dessemelhanca entre os effeitos e a causa. Por isso, \emph{todo} o effeito tem com a sua causa uma proporção de \emph{ordem}, emquanto é produzido por ella; mas nem \emph{todo} o effeito tem com a causa uma proporção de \emph{natureza}. Appliquemos estes principios ao nosso caso. Os seres visiveis, sendo finitos e contingentes, não exprimem toda a virtude da sua causa, não possuem á mesma \emph{natureza} de Deus, e por isso não nos levam ao conhecimento \emph{adequado} da \emph{natureza} de Deus. Mas, sendo \emph{effeitos}, e effeitos que só podem derivar de um Ente necessario e omnipotente, levam certamente ao conhecimento da \emph{existencia} d'esse Ente. --- Todavia, embora o conhecimento, que temos da \emph{natureza} de \textsc{Deus}, não seja \emph{adequado}, mas \emph{imperfeito}, comtudo é \emph{positivo}; pois as coisas creadas são semelhanças \emph{positivas}, embora \emph{imperfeitas}, das divinas perfeições, visto que todo o agente tende a imprimir no effeito a propria \emph{semelhança}.}.

%%>>> art-37

\section*{Artigo II.\\ Demonstração da existencia de Deus.}
\addcontentsline{toc}{section}{Artigo II. --- Demonstração da existencia de Deus.}

11. \textbf{Especies de argumentos para} a \textbf{demonstração da existencia de Deus}. --- A existencia de \textsc{Deus} demonstra-se \emph{a-posteriori}. Ora os argumentos, que demonstram \emph{a-posteriori} a existencia de \pg{427} \pg{428} \textsc{Deus}, são de tres especies: \emph{metaphysicos}, \emph{physicos} e \emph{moraes}. Os \emph{metaphysicos} basé6am-se na propria \emph{natureza} das coisas, --- os \emph{physicos}, na \emph{ordem} que existe no mundo e, em particular, nas 